\chapter*{Conferencia XII}
\addcontentsline{toc}{chapter}{LECTURA I - El timbre y el carácter del violín}
SEÑORES: (g) La posición del puente es un modificador importante del tono. Es tan importante que una mala posición puede aniquilar por completo la "riqueza" del sonido del violín, incluso cuando todos los demás factores están en su mejor momento. Recuerdo ese caso de "wolf" que se curó completamente con la posición del puente y el poste. También recuerdo que los armónicos a bassa, o tonos resultantes de los libros de texto, tonos que ayudan a producir un sonido de violín "rico", solo se harán audibles cuando la longitud de las cuerdas (desde el arco hasta la cejuela, no desde el puente hasta la cejuela) y la longitud activa de la caja de resonancia sean iguales. Por lo tanto: Debido a que el arco prácticamente acorta las cuerdas en una pulgada, y debido a que la mayor longitud activa de la caja de resonancia, que produce sonido audible, es igual a 12 pulgadas, por lo tanto, para la producción del violín "rico"
tono, la longitud de las cuerdas, desde el puente hasta '
tuerca,
Debe medir 13 pulgadas. Sin embargo, en la regulación del tono, esta regla no debe interpretarse como una norma estricta. Según mi experiencia, existe un único punto para la posición del puente que produce la mayor riqueza tonal; y, de ninguna manera, puedo predeterminar con precisión dicha posición. Para determinar con exactitud la mejor posición del puente, es necesario experimentar con cada caja de resonancia.
Concluyo la discusión sobre el puente del violín con la afirmación categórica de que todos los puentes de madera inmadura (siempre albura) se desechan mejor como
leña.
(8)	El problema del diapasón viene a continuación en el 180.

Lista de factores que influyen en la disminución del volumen y la intensidad del sonido del violín. Debido a mi incapacidad para resolver el problema del puente hasta después de resolver el del diapasón, este último recibió prioridad en la presentación. Por el momento, no puedo añadir nada más sobre este modificador del sonido.
(9) Las cuerdas son un factor muy importante en la disminución de la potencia sonora del violín. Su importancia se evidencia claramente en el hecho de que la rigidez de la caja de resonancia debe estar determinada por el diámetro y el peso de las cuerdas. En mi declaración sobre las cuerdas de violín, la cuerda de alambre no recibe más consideración que la de un mal necesario debido a situaciones adversas; es decir, situaciones en las que el exceso de vapor de agua provoca el rápido deterioro de las cuerdas de tripa. En todos los casos en que la rigidez de la caja de resonancia está determinada por cuerdas gruesas, la sustitución por cuerdas más finas disminuye la potencia sonora. He observado que en todos los casos en que la rigidez de la caja de resonancia se reduce precisamente para corresponder con la fuerza en cuerdas de calibre 2, la sustitución por cuerdas más gruesas aumenta el volumen del sonido a la vez que disminuye su intensidad. Por lo tanto, a corta distancia el sonido es mayor, pero falla a larga distancia. Considero que una explicación satisfactoria para este fenómeno es compleja. Como ayuda para dicha explicación, señalo que las cuerdas más gruesas, tocadas en dobles cuerdas, provocan una vibración violenta de la caja de resonancia. Parece razonable suponer que el temblor del violín puede ser tan grande como para debilitarlo.
181

tono en vista del hecho de que el temblor violento de la trompa opera no solo debilita el tono, sino también, perceptiblemente, baja la altura del tono. Es un hecho de observación más o menos frecuente que el tono forzado de cualquier trompa en una banda de metales está desafinado con respecto a las demás trompas, y desafinado debido a que su altura del tono ha disminuido. Me parece que dicha altura del tono es causada por el temblor violento del instrumento. Que dicho tono se debilite en intensidad es claramente perceptible. Asimismo, en un conjunto orquestal, el tono del violín que tiembla violentamente puede ser totalmente aniquilado por las ondas sonoras de los instrumentos armónicos. Para este hecho, no parece haber una razón más plausible que la disminución de la intensidad en el tono de dicho violín tembloroso. De nuevo, en el coro, cuando algún miembro se entrega a un trémolo violento, es inmediatamente evidente que su voz (demasiado a menudo "su", ¡qué lástima!) no solo se debilita, sino que está tan desafinada que resulta desagradable para el oyente con formación musical. Nuevamente, las cuerdas de diámetro demasiado grande tienden a acentuar la calidad del tono ruidoso. Por lo tanto, cualquiera que sea la explicación correcta, el hecho es que el uso de cuerdas con demasiada fuerza para la rigidez de la caja de resonancia, o para la rigidez de la placa posterior, es un desastre asegurado. Desde mi punto de vista, es más seguro errar en la dirección opuesta porque las cuerdas más pequeñas tienden a disminuir el ruido.
calidad del tono.
El "giro" y el número de "hebras" en el
Las cuerdas pueden funcionar para disminuir tanto el volumen como
intensidad del tono.	Por lo tanto, la cadena de un solo 182

La flexibilidad de las cuerdas disminuye la potencia del sonido. La razón es obvia. En este tipo de cuerdas, al ser mínima la flexibilidad, la rapidez para enrollarlas y desenrollarlas bajo la presión del arco disminuye considerablemente; por lo tanto, la fuerza de los golpes sobre el puente se reduce. La flexibilidad de las cuerdas de violín aumenta la rapidez para enrollarlas y desenrollarlas; por consiguiente, la cuerda flexible produce un golpe de mayor fuerza. El peso de las cuerdas también puede disminuir la potencia del sonido. Para que tengan el mismo peso, la cuerda de seda debe tener un diámetro mayor, mientras que la de acero debe tener un diámetro menor. En todos los casos, la uniformidad del sonido se ve afectada por la desproporción entre el diámetro y el peso de las cuerdas. Como modificador del sonido del violín, el tema de las cuerdas despierta gran interés entre los estudiosos de las peculiaridades del sonido. El sonido del mejor violín jamás construido puede arruinarse fácilmente por las cuerdas utilizadas. En mi opinión, la elección del metal para enrollar la cuerda de sol debe regirse por las peculiaridades del sonido inherentes a cada violín. En igual cantidad, el cobre, la plata y el oro varían en peso. Siempre que he encontrado una cuerda entorchada de cobre que se ajusta con precisión a las peculiaridades tonales de un violín determinado, la sustitución por una cuerda G entorchada de plata u oro ha resultado ser una decepción.
Para un ajuste preciso de la cuerda G, mi método consiste en probar cuerdas de diferentes diámetros. Después de determinar el diámetro que da mejores resultados en los acordes dobles, selecciono una de diámetro ligeramente mayor y la dejo estirar. Luego, estando en posición y afinada,
183

Proceda a disminuir tanto el peso como el diámetro con un trozo de tela de esmeril de una pulgada cuadrada. Con precisión, la tela de esmeril, doblada una vez alrededor de la cuerda y sujetada firmemente con el pulgar y el índice, se desplaza desde el puente hasta la cejuela, luego se gira ligeramente y se vuelve a colocar en el puente. Esta operación se interrumpe frecuentemente al aplicar el arco, para así determinar el momento en que la reducción de diámetro y peso ha alcanzado el punto deseado. Si se realiza con cuidado, este método elimina en gran medida ciertas irregularidades tonales inevitables en todas las cuerdas de sol nuevas. Si se aplica a una cuerda de sol con el diámetro correcto, el resultado es desastroso. Considerando el gran valor de una hermosa calidad de sonido en una cuerda de sol, ningún esfuerzo es excesivo.
[Es posible que me acusen de repetición frecuente; pero asumiré ese riesgo al afirmar nuevamente que los hermosos tonos de doble cuerda son la salvación del violinista solista. Con la salvación en mente, la cuerda de sol no debe ser ofensivamente prominente. Es posible que sea difícil de complacer en lo que respecta a la calidad del sonido del violín; pero, todo lo que diga sobre este punto es solo mi opinión personal, y la digo reconociendo que mi opinión puede ser errónea. En la vida, no conozco un hecho más poderoso que el hecho de que el gusto musical varía según el número de personas. En mi opinión, el sonido de un juego de cuerdas de acero para violín es digno de anatema.
(10) Que el estado de las superficies interiores del violín
184

Puede funcionar para disminuir tanto el volumen como la intensidad del tono, solo requiere un pensamiento. ¿Recomendé extender una alfombra sobre la superficie interior del?
Tras el violín, sin duda cualquier lector me juzgaría así: «Ese Castle es un tonto».
o peor aún, "Castle es un idiota". Lo peor de todo es que "el abrigo le quedaría bien". Sin duda, muchos lectores piensan que "el abrigo les queda bien" cuando digo que tanto el volumen como la intensidad del sonido del violín se ven aumentados por una superficie interior perfectamente lisa. Lo digo, y lo digo con toda la sinceridad de la que soy capaz. Este problema en el sonido del violín está resuelto a mi satisfacción. Si alguna buena calidad de sonido se viera perjudicada por la superficie interior perfectamente lisa, entonces de ninguna manera este problema estaría resuelto a mi satisfacción. La cuestión de resolver este problema a su satisfacción recae enteramente en ustedes. Jamás les daría los detalles del trabajo de la superficie interior si "solo estuviera buscando salud". Al contrario, me guardaría para mí el secreto de convertir violines de 5 dólares en violines de 100 dólares. Si es necesario, puedo darles los nombres de varios violinistas bastante competentes cuyos violines de 3 dólares (al por mayor) ahora son valorados por ellos en varios cientos de dólares. No me entiendan como si afirmara que todo este cambio en el valor del sonido se debe a una superficie interior perfectamente lisa; sino que pueden entender que digo que, sin una superficie interior permanentemente lisa, Los propietarios no habrían fijado esos precios tan altos. Puedo dar el nombre de uno de esos propietarios, que vive a 30 minutos en coche de uno de los mercados de violines más grandes del mundo,
185

quien se negó rotundamente a ponerle precio a su
Originalmente \$10-flddle.
¿Les di detalles a esos dueños? —Por supuesto que no. —¿Por qué no? —Precisamente por el prejuicio que usted alberga en este momento.
Aquí hay un "consejo" destinado a su beneficio: Cuando alguien garantiza que el tono de su violín "se propagará" a 1250 pies medidos al aire libre y bajo condiciones meteorológicas idénticas, y en las mismas horas del día, y con las mismas precauciones contra la presencia de reflectores de ondas sonoras que se han descrito hasta ahora, simplemente tómese la molestia de determinar si la superficie interior de ese violín es permanente y perfectamente
liso.
De nuevo, en todos los violines antiguos que no han sido "limpiados", y en muchos violines usados que aún no son viejos, les aseguro que encontrarán su superficie interior cubierta con una "alfombra". Dicha alfombra estará compuesta de fibras de madera, polvo de madera y suciedad. Esta es su oportunidad de comprobar el aumento de volumen e intensidad del sonido sin dañar la calidad del tono, de la cual he dado garantías en estas páginas. A todos los violinistas que deseen un violín con la máxima potencia sonora, les recomiendo la aplicación cuidadosa de las instrucciones para la protección de la superficie interior*, junto con las demás instrucciones aquí presentadas. Con un trabajo cuidadoso y minucioso, estoy seguro de que tendrán éxito.
(11) El área y la posición de las salidas son los más
186

Existen factores poderosos que operan para disminuir el volumen y la intensidad del sonido del violín. El área de salidas, por sí sola, es más potente que todos los demás factores combinados. Por sorprendente que parezca esta afirmación, su veracidad es fácil de demostrar. Sin salidas, los mejores violines jamás construidos no producirán, ni podrán producir, más sonido que un violín de palo de escoba. En tal situación, las cuerdas se quedan solo con aire no confinado sobre el cual ejercer su fuerza. Por lo tanto, en lugar de una concentración de las líneas de propagación de las ondas sonoras, se produce la mayor dispersión posible de dichas líneas. Además, en tal situación, la fuerza de las cuerdas no recibe ningún aumento ni de la acción directa de la caja de resonancia ni de la acción simpática. Sin salidas, la caja de resonancia está inmóvil. Su potencia no es suficiente para comprimir el aire confinado. Sin salidas, el violín
desciende al nivel del palo de escoba. Por lo tanto, sin salidas sin violín.
Pero una pequeña abertura en las paredes permite cierta acción de la caja de resonancia. Inmediatamente se produce un aumento perceptible de la potencia del sonido; y dicho aumento se produce al colocar la abertura en la parte posterior, en las costillas, en la caja de resonancia, donde se desee, pero dicho aumento no se produce en igual medida. Esas paredes curvas dirigen las líneas de propagación de las ondas sonoras. Al observar un violín de buena calidad, se percibe inmediatamente que esas paredes curvas internas no pueden dirigir la bola que rebota hacia las costillas. Si esas costillas fueran de cristal, ninguna bola que rebota las rompería.
187

La pelota que rebota va, allí también, la onda sonora.
El movimiento es fundamental. Sin embargo, existe una diferencia entre la acción de una sola bola y la de innumerables bolas que se tocan entre sí. Las moléculas de aire, al ser bolas que poseen una inmensa energía elástica y no moverse de la posición que ocupan al recibir un golpe, producen ondas sonoras al transmitir energía elástica de una a otra; y dicha transmisión continúa hasta que la fuerza del golpe original se agota. Es evidente que aumentar la fuerza del golpe original incrementa la distancia recorrida por la onda sonora. También es evidente que el movimiento de la onda sonora puede dispersarse o concentrarse; además, que la concentración de dicho movimiento aumenta tanto el volumen como la intensidad del tono, pero en grados desiguales; la intensidad aumenta en mayor medida. Es evidente que las paredes internas del cuerpo del violín no solo impiden la dispersión del movimiento de la onda sonora, sino que, debido a la curvatura longitudinal y transversal de las tapas, también lo concentran. Es evidente que el punto de concentración de la onda sonora dentro del violín debe variar con los diferentes grados de curvatura de las tapas. Determinar con precisión estos puntos variables es el gran desafío de los científicos. Es evidente que ubicar las salidas en los puntos de mayor concentración de ondas sonoras se convierte en un factor importante para la producción de la máxima potencia tonal; asimismo, ubicar las salidas a una distancia de dichos puntos de concentración reduce la potencia tonal.
188

¿Cómo podemos encontrar esos puntos? No le preguntes al científico.
Él no puede decirte más de lo que Sam Jones puede decirte dónde está el cielo.
Basándose únicamente en la experiencia, el constructor de violines tallará las salidas oblicuamente a través de la caja de resonancia y confiará en la suerte para acertar en esos puntos. De este modo, la oblicuidad en la posición de las salidas adquiere un valor inmenso para la potencia del sonido del violín. Si se pudieran predeterminar los puntos de mayor concentración de ondas sonoras, se podrían emplear salidas en forma de paralelogramo para aumentar la potencia del sonido.
Como se ha demostrado anteriormente, la salida de un área pequeña permite un grado limitado de acción de la caja de resonancia; por lo tanto, se deduce que aumentar el área de salidas permite una mayor amplitud de oscilación de la caja de resonancia; por consiguiente, una mayor fuerza de soplado sobre el aire contenido; por consiguiente, una mayor potencia tonal.
Hemos llegado a un punto en la construcción del violín que posee un intenso interés para dos clases de usuarios de violín; uno que desea calidad de tono con volumen moderado; el otro que desea gran volumen de tono independientemente de la calidad: Desde mi punto de vista, el interés en las salidas de violín es igual al interés en
Madera de caja de resonancia y modelo de arco de placa. El análisis de las salidas continuará en la próxima hora.
